\documentclass{article}
\begin{document}
\tableofcontents
\section{Scope}
The record defines the material covered by this catalogue.
\section{Sources}
The field notes and specimens come from two separate surveys.
\section{Dates}
Each date identifies the observation rather than the publication.
\section{Methods}
Measurements follow a consistent unit and ordering rule.
\section{Preparation}
The samples were prepared before any comparison was made.
\section{Collection}
Each specimen has a collection site and a short description.
\section{Weather}
Conditions at the site affected the initial observations.
\section{Storage}
The storage record preserves the order of collection.
\section{Inspection}
The inspection compares visible features without inference.
\section{Comparison}
The comparison uses the same criteria for every specimen.
\section{Exceptions}
Unexpected measurements are recorded without correction.
\section{Review}
A separate review checks the descriptions against the notes.
\section{Provenance}
The provenance record traces each item to its field entry.
\section{Limitations}
The catalogue excludes items without a confirmed origin.
\section{Summary}
The summary lists the supported observations only.
\section{Archive}
The archive retains the source notes for later inspection.
\end{document}
